% !PRINTED 508
\weberpage{508}

Hieran schliesst sich der \gesperrt{Satz von Lagrange}\footnote{\gesperrt{Lagrange}, Réflexions sur la résolution algébrique des équations. Mémoires de l'Académie de Berlin, années 1770, 1771. Oeuvres de Lagrange. Tome III. Der Satz ist von \gesperrt{Lagrange} allerdings nur in einer specielleren Fassung gegeben. Die allgemeine Formulirung rührt von \gesperrt{Galois} her.}:

2. \gesperrt{Jede Grösse des Körpers $N$, die die Permutationen
 der Gruppe $Q$ gestattet, ist in dem
Körper $\Omega(\psi)$ enthalten, wenn $\psi$ eine zu $Q$ gehörige
 Function ist.}

Eine die Permutationen von $Q$ gestattende Function $\omega$ geht
durch die Permutationen einer Nebengruppe $Q\pi_1$ in ein und
dieselbe Function über. Es entsprechen also den conjugirten
Werthen

\weq{3}{\psi, \psi_1, \psi_2 \ldots \psi_{j-1}}

die Werthe

\weq{4}{\omega, \omega_1, \omega_2 \ldots \omega_{j-1},}

die jedoch nicht nothwendig alle von einander verschieden sind.

Wendet man auf die Grössenreihen (3), (4) eine der Permutationen
 von $P$ an, so tritt eine gewisse Permutation ein, und
zwar in beiden Reihen die gleiche, da, wenn z. B. $\psi_1$ in $\psi_2$ übergeht,
auch $\omega_1$ in $\omega_2$ übergehen muss (§. 154, 5.).

Betrachten wir nun die Summe

\[
\varphi(t) \left( \frac{\omega}{t - \psi} + \frac{\omega_1}{t - \psi_1} + \cdots + \frac{\omega_{j-1}}{t - \psi_{j-1}} \right) = \chi(t),
\]

die eine ganze Function $(j-1)^{\text{ten}}$ Grades von $t$ ist, so finden
wir, dass sie durch alle Permutationen $P$ ungeändert bleibt und
folglich in $\Omega$ enthalten ist. Setzt man dann $t = \psi$ und beachtet,
dass $\varphi(t)$ keine gleichen Wurzeln hat, so folgt

\weq{5}{\omega = \frac{\chi(\psi)}{\varphi'(\psi)}.}

Es ist also $\omega$ rational durch $\psi$ ausgedrückt, und dies ist der
Inhalt des Satzes 2.

3. \gesperrt{Wenn wir eine zu $Q$ gehörige Function $\psi$ dem
Körper $\Omega$ adjungiren, so reducirt sich die
Gruppe des Körpers $N$ auf $Q$.}

% [page alignment corrected 2026-08-07] Printed page 508 continues through the
% line ending "weil $\psi$ durch".  The following page begins with "$Q$
% ungeändert bleibt"; the displaced line was moved here from i0534.
Denn bezeichnen wir den Körper $\Omega(\psi)$ mit $\Omega'$, so gestattet
\gesperrt{erstens} jede Gleichung in $\Omega'$ zwischen $\alpha, \alpha_1 \ldots \alpha_{m-1}$ die Permutationen
von $Q$, weil $Q$ in $P$ enthalten ist, und weil $\psi$ durch
%CONT>>
